% ---------------------------------------------------------------------
%  PHỤ LỤC D. MA TRẬN TRUY VẾT YÊU CẦU
% ---------------------------------------------------------------------
\section{Ma trận truy vết yêu cầu}\label{pl:truy-vet}

Ma trận truy vết tại Bảng~\ref{tab:truy-vet} liên kết yêu cầu trong đề bài~\cite{debai} và chức năng mở rộng ở Bảng~\ref{tab:chuc-nang-mo-rong} với nghiệp vụ con, use case, quy tắc nghiệp vụ và yêu cầu phi chức năng tương ứng. Yêu cầu đề bài được mã hóa theo thứ tự xuất hiện (YC); ba đặc điểm nổi bật được ký hiệu ĐĐ. Ma trận hỗ trợ kiểm tra theo hai chiều: mỗi yêu cầu phải được hiện thực bởi ít nhất một use case và mỗi use case phải truy ngược được về ít nhất một yêu cầu.

{\small
\begin{longtable}{|C{1.2cm}|L{\dimexpr\textwidth-12\tabcolsep-7\arrayrulewidth-1.2cm-2.0cm-2.6cm-1.9cm-1.7cm\relax}|L{2.0cm}|L{2.6cm}|L{1.9cm}|L{1.7cm}|}
\caption{Ma trận truy vết từ yêu cầu đến nghiệp vụ và use case}\label{tab:truy-vet}\\
\hline
\hd{Mã} & \hd{Yêu cầu} & \hd{Nghiệp vụ} & \hd{Use case} & \hd{Quy tắc} & \hd{NFR}\\ \hline
\endfirsthead
\multicolumn{6}{l}{\footnotesize\itshape Bảng \thetable\ (tiếp theo)}\\ \hline
\hd{Mã} & \hd{Yêu cầu} & \hd{Nghiệp vụ} & \hd{Use case} & \hd{Quy tắc} & \hd{NFR}\\ \hline
\endhead
\multicolumn{6}{|l|}{\hs{Yêu cầu của đề bài}}\\ \hline
YC00 & Định danh người dùng làm nền tảng cho chia sẻ và phân quyền (yêu cầu ngầm định) & NV01.1–NV01.4 & UC01–UC07 & BR01–BR03 & NFR14, NFR15, NFR17\\ \hline
YC01 & Tạo danh sách mua sắm theo từng ngày, chỉnh sửa, cập nhật và lưu lại cho những lần sau & NV02.1–NV02.3 & UC16, UC17, UC18, UC23 & BR08, BR09 & NFR05, NFR11\\ \hline
YC02 & Tạo nhóm gia đình, chia sẻ danh sách, mọi thành viên cùng đóng góp và cập nhật & NV01.5–NV01.7, NV02.2 & UC09–UC15, UC17 & BR04–BR07, BR23 & NFR10, NFR16\\ \hline
YC03 & Trưởng nhóm điều phối, phân công mua sắm, bảo đảm không quên hoặc trùng lặp & NV02.4, NV02.5 & UC19, UC20 & BR09, BR10 & NFR10\\ \hline
YC04 & Lên kế hoạch bữa sáng, trưa, tối theo ngày hoặc tuần & NV05.1, NV05.2 & UC37, UC38, UC39 & BR18 & NFR05\\ \hline
YC05 & Đề xuất món ăn dựa trên thực phẩm có sẵn trong tủ lạnh & NV05.3 & UC40 & BR19 & NFR09\\ \hline
YC06 & Lưu thông tin thực phẩm: tên, số lượng, ngày hết hạn, vị trí lưu trữ; theo dõi và quản lý & NV03.1–NV03.4 & UC24, UC25, UC27, UC28, UC29 & BR11, BR12, BR14, BR15 & NFR04\\ \hline
YC07 & Tự động nhắc khi thực phẩm sắp hết hạn trước 3 ngày & NV03.5 & UC08, UC30, UC31 & BR12, BR13 & NFR12\\ \hline
YC08 & Lưu, quản lý công thức yêu thích, liên kết nguyên liệu với tủ lạnh & NV04.1, NV04.2, NV04.4 & UC32, UC33, UC34, UC35 & BR16, BR17 & —\\ \hline
YC09 & Tìm kiếm món ăn, thực phẩm trong hệ thống hoặc trong tủ lạnh & NV06.1 & UC42 & BR23 & NFR09\\ \hline
YC10 & Phân loại thực phẩm theo danh mục như thịt, rau củ, đồ hộp & NV06.2, NV08.3 & UC43, UC50 & BR22 & —\\ \hline
YC11 & Gợi ý món ăn phù hợp với thực phẩm sẵn có & NV06.3 & UC30, UC40 & BR19 & NFR09\\ \hline
YC12 & Báo cáo đồ đã mua và thực phẩm đã tiêu thụ trong một khoảng thời gian & NV07.1, NV07.2 & UC44, UC45 & BR21 & —\\ \hline
YC13 & Quản trị viên quản lý tài khoản người dùng & NV08.1 & UC48 & BR23 & NFR16\\ \hline
YC14 & Quản trị viên quản lý danh mục tên món đồ, loại mặt hàng, đơn vị đo và cấu hình hệ thống & NV08.2–NV08.5 & UC49–UC53 & BR16, BR22 & NFR20\\ \hline
\multicolumn{6}{|l|}{\hs{Đặc điểm nổi bật của đề bài}}\\ \hline
ĐĐ01 & Tính tiện lợi cao, kiểm soát toàn bộ quy trình mua sắm và nấu nướng & NV01–NV08 & Toàn bộ & — & NFR05, NFR08, NFR11\\ \hline
ĐĐ02 & Chia sẻ trong gia đình, phân công giữa các thành viên & NV01.5–NV01.7, NV02.4 & UC09–UC15, UC19 & BR04–BR07, BR10 & NFR10, NFR16\\ \hline
ĐĐ03 & Thông báo nhắc nhở thông minh, chọn món từ nguyên liệu sẵn có & NV03.5, NV05.3 & UC08, UC31, UC40 & BR13, BR19 & NFR12\\ \hline
\multicolumn{6}{|l|}{\hs{Chức năng mở rộng đề xuất}}\\ \hline
MR01 & Mời thành viên bằng mã mời hoặc mã QR & NV01.6 & UC10, UC11 & BR06 & NFR17\\ \hline
MR02 & Chuyển mặt hàng đã mua vào tủ lạnh & NV02.6 & UC21 & BR11 & NFR13\\ \hline
MR03 & Sinh danh sách mua sắm từ kế hoạch bữa ăn & NV02.7 & UC22 & BR20 & —\\ \hline
MR04 & Quét mã vạch sản phẩm & NV03.1 & UC26 & BR22 & NFR18\\ \hline
MR05 & Gợi ý món có chấm điểm, ưu tiên thực phẩm sắp hết hạn & NV05.3 & UC40 & BR19 & NFR09\\ \hline
MR06 & Xác nhận đã nấu, trừ kho theo FEFO & NV05.4 & UC41 & BR15 & NFR13\\ \hline
MR07 & Bổ sung nguyên liệu còn thiếu từ công thức & NV04.3 & UC36 & BR09 & —\\ \hline
MR08 & Báo cáo chi tiêu, lãng phí, xuất báo cáo & NV07.1, NV07.3, NV07.4 & UC44, UC46, UC47 & BR21 & —\\ \hline
MR09 & Thiết lập nhắc nhở cá nhân hóa & NV01.4 & UC08 & BR12, BR13 & —\\ \hline
MR10 & Đồng bộ thời gian thực và hoạt động ngoại tuyến & NV02.2, NV02.5 & UC16, UC17, UC20 & — & NFR10, NFR11\\ \hline
\end{longtable}
}

Kiểm tra ngược cho thấy cả 53 use case đều có trong ma trận. Trong số đó, 46 use case hiện thực các yêu cầu YC00–YC14; bảy use case còn lại — \uc{UC21}, \uc{UC22}, \uc{UC26}, \uc{UC36}, \uc{UC41}, \uc{UC46} và \uc{UC47} — phát sinh từ chức năng mở rộng. Các use case \uc{UC08}, \uc{UC10}, \uc{UC11}, \uc{UC40} và \uc{UC44} vừa đáp ứng yêu cầu đề bài, vừa được mở rộng tương ứng. Kết quả truy vết không ghi nhận use case nằm ngoài phạm vi hoặc yêu cầu nào chưa được hiện thực.
